\documentclass[Screen16to9,17pt]{foils}
\usepackage{zencurity-slides}

\selectlanguage{english}

\usepackage{tabularx}

% Reused from pentest-II-foredrag + hackerworkshop-prosa (obj 7).

\begin{document}

\mytitlepage{Web Attacks I -- Attacks}
{Cyber Akademiet, 2026}

\slide{Goals for today}
\begin{list2}
\item Identify and explain common web-application weaknesses
\item SQL injection, cross-site scripting, cross-site request forgery
\item Session handling, file inclusion and security misconfiguration
\item Relate them to the OWASP Top 10
\end{list2}

\slide{Plan for today}
\begin{list2}
\item Why web hacking; the OWASP Top 10 as a map
\item Injection: SQL injection
\item Client-side: XSS and CSRF
\item Sessions, TLS/Heartbleed and web scanning
\end{list2}

\slide{Time schedule}
\begin{list2}
\item 08:30--10:00 Web pentest intro; OWASP Top 10; SQL injection
\item 10:00--10:15 Break
\item 10:15--11:45 XSS, CSRF, sessions, TLS/Heartbleed
\item 11:45--12:30 \textbf{Exercise} (booklet)
\item 12:30--13:15 Lunch
\end{list2}

\slide{Why talk about hacking web}

\hlkimage{7cm}{pawel-janiak-dxFi8Ea670E-unsplash.jpg}

\begin{list2}
\item Everything is web today
\item HTTPS/TLS used for data transport in most applications
\item Web APIs are used across the industry
\item Even \emph{desktop} and \emph{mobile} application are rarely more than a specialized browser\\ like \url{electronjs.org} -- Chromium and Node.js
\hfill {\footnotesize Photo by Pawel Janiak on Unsplash}
\item Also software security and web security has a HUGE overlap
\end{list2}

\slide{Web Application Security}

\hlkimage{6cm}{hoffman-web-application-security.jpg}
\emph{Web Application Security}, Andrew Hoffman, 2020, ISBN: 9781492053118 called WAS

\slide{Primary HTTP mthods}

\begin{list2}
\item [GET]
Requests a representation of the specified resource. Requests using GET should only retrieve data and should have no other effect. (This is also true of some other HTTP methods.)[1] The W3C has published guidance principles on this distinction, saying, "Web application design should be informed by the above principles, but also by the relevant limitations."[13] See safe methods below.
\item [POST]
Requests that the server accept the entity enclosed in the request as a new subordinate of the web resource identified by the URI. The data POSTed might be, for example, an annotation for existing resources; a message for a bulletin board, newsgroup, mailing list, or comment thread; a block of data that is the result of submitting a web form to a data-handling process; or an item to add to a database.[14]
\item [PUT]
Requests that the enclosed entity be stored under the supplied URI. If the URI refers to an already existing resource, it is modified; if the URI does not point to an existing resource, then the server can create the resource with that URI.[15]
\end{list2}

Source: \link{https://en.wikipedia.org/wiki/Hypertext_Transfer_Protocol}

\slide{OWASP Top Ten}

\hlkimage{16cm}{owasp.jpg}

\begin{quote}
The OWASP Top Ten provides a minimum standard for web application
security. The OWASP Top Ten represents a broad consensus about what
the most critical web application security flaws are.
\end{quote}

\begin{list1}
\item The Open Web Application Security Project (OWASP)
\item OWASP produces lists of the most common types of errors in web applications
\item \link{http://www.owasp.org}
\item Create Secure Software Development Lifecycle
\end{list1}

\slide{SQL injection}

SQL Injection FAQ from Chip Andrews 2001 \link{http://www.sqlsecurity.com}:

\begin{minted}[fontsize=\footnotesize]{shell}
Set myRecordset = myConnection.execute
("SELECT * FROM myTable
WHERE someText ='" & request.form("inputdata") & "'")
\end{minted}

Using input \verb+exec master..xp_cmdshell 'net user test testpass /ADD' --+

Will perform
\begin{minted}[fontsize=\footnotesize]{shell}
SELECT * FROM myTable
WHERE someText ='' exec master..xp_cmdshell 'net user test testpass /ADD'--
\end{minted}

Depending on the access rights this might have full authorization in the database!

\slide{Sqlmap}

\begin{quote}\footnotesize
sqlmap is an open source penetration testing tool that automates the process of detecting and exploiting SQL injection flaws and taking over of database servers. It comes with a powerful detection engine, many niche features for the ultimate penetration tester and a broad range of switches lasting from database fingerprinting, over data fetching from the database, to accessing the underlying file system and executing commands on the operating system via out-of-band connections.

\end{quote}

\begin{list1}
\item Automatic SQL injection and database takeover tool
\link{http://sqlmap.org/}
\end{list1}

\slide{sqlmap features}

\hlkimage{15cm}{sqlmap-features-1.png}

Not a complete list!

Source: \link{http://sqlmap.org/}

\slide{10. Cross-Site Scripting (XSS)}

\begin{quote}\small
  Cross-Site Scripting (XSS) vulnerabilities are some of the most common vulnerabili‐
  ties throughout the internet, and have appeared as a direct response to the increasing
  amount of user interaction in today’s web applications.
  At its core, an XSS attack functions by taking advantage of the fact that web applica‐
  tions execute scripts on users’ browsers. Any type of dynamically created script that is
  executed puts a web application at risk if the script being executed can be contamina‐
  ted or modified in any way—in particular by an end user.
\end{quote}
Source: \emph{Web Application Security}, Andrew Hoffman, 2020, ISBN: 9781492053118

XSS attacks are categorized a number of ways, with the big three being:
\begin{list2}
\item Stored (the code is stored on a database prior to execution)
\item Reflected (the code is not stored in a database, but reflected by a server)
\item DOM-based (code is both stored and executed in the browser)
\end{list2}

Note: attacks the user, his/her data mostly.

\slide{Cross-site scripting}

\begin{list1}
\item When logged into a site, we have a session identifier -- being presented with each HTTP request
\item If an attacker can \emph{activate} javascript that access and re-send this to another site we have a cross-site scripting attack
\item We often check using some parameter with code, like this:
\begin{alltt}\footnotesize
<A HREF="http://example.com/comment.cgi?
mycomment=<SCRIPT>malicious code</SCRIPT>
">Click here</A>
\end{alltt}
\item If this code is returned, and active, as part of the returned page and HTML we have a cross-site scripting possibility
\item This code is executed in the user browser, with the same permissions as the user
\end{list1}

\slide{11. Cross-Site Request Forgery (CSRF)}

\begin{quote}
Sometimes we already know an API endpoint exists that would allow us to perform an operation we wish to perform, but we do not have access to that endpoint because it requires privileged access (e.g., an admin account).

In this chapter, we will be building Cross-Site Request Forgery (CSRF) exploits that result in an admin or privileged account performing an operation on our behalf rather than using a JavaScript code snippet.

CSRF attacks take advantage of the way browsers operate and the trust relationship between a website and the browser. By finding API calls that rely on this relationship to ensure security—but yield too much trust to the browser—we can craft links and forms that with a little bit of effort can cause a user to make requests on his or her own behalf—unknown to the user generating the request.
\end{quote}

\begin{list2}
  \item Often seen in small CPE routers in homes, if the user activates and evil link, their router might be reconfigured or taken over
\end{list2}

\slide{Session IDs}

\begin{list2}
\item Session IDs tie the user with the state on the server
\item Must be randomly assigned, otherwise an attacker can guess a valid ID
\item Common problems, time based or predictable in some way
\item Check code for generating IDs or measure - Phase Space Analysis
\end{list2}

\slide{Testing Web APIs 2024}

\verb+https://customerdomain.example.com/api/customer/v1/business/$ENDPPOINT+\\
ENDPOINT being something like: employeedocuments, clientcontacts, businessPropertyObjects, activities

\begin{list2}
\item    We love REST and the world does too
\end{list2}

\slide{Generic Network Fault Injection}

\begin{list1}
\item Inserting proxies can allow modification of data in transit
\item Can be used for random bit corruption
\item Can often reproduce the data
\item Automate gathering of evidence
\item You can use a simple Random TCP/UDP fault injector, but usually you would get better results by conforming to some structure
\item Various test cases must tried with potential bad data, examples:
\begin{list2}
\item loooong input in all fields and input -- buffer overflows
\item SQL injection -- trying database commands as input
\item Cross-site scripting -- does inserting JavaScript result in this being included in pages returned
\item Random bytes - recommend using real fuzzers that understand target protocol
\item Metacharacters like null bytes, or \verb+.,:;/\textbackslash+
\end{list2}
\end{list1}

\slide{Transport Layer Security (TLS)}

\hlkimage{5cm}{crypto-class.png}

\begin{list1}
\item Originally from Netscape Communications Inc.
\item Secure Sockets Layer (SSL) was adopted by the IETF newer versions are called
Transport Layer Security (TLS)
\item TLS 1.0 based on generalized version of SSL Version 3.0
\item RFC-2246 The TLS Protocol Version 1.0 from Januar 1999
\item RFC-3207 SMTP STARTTLS allows the use of TLS with SMTP mail protocols
\item Today most sites and servers support TLS Server Name Indication (SNI) name based certificates
\end{list1}

\slide{Heartbleed CVE-2014-0160}

\hlkimage{18cm}{heartbleed-com.png}

\centerline{Nok det mest kendte SSL/TLS exploit}

Source: \link{http://heartbleed.com/}

\slide{Heartbleed hacking}

\begin{alltt}\footnotesize
  06b0: 2D 63 61 63 68 65 0D 0A 43 61 63 68 65 2D 43 6F  -cache..Cache-Co
  06c0: 6E 74 72 6F 6C 3A 20 6E 6F 2D 63 61 63 68 65 0D  ntrol: no-cache.
  06d0: 0A 0D 0A 61 63 74 69 6F 6E 3D 67 63 5F 69 6E 73  ...action=gc_ins
  06e0: 65 72 74 5F 6F 72 64 65 72 26 62 69 6C 6C 6E 6F  ert_order&billno
  06f0: 3D 50 5A 4B 31 31 30 31 26 70 61 79 6D 65 6E 74  =PZK1101&payment
  0700: 5F 69 64 3D 31 26 63 61 72 64 5F 6E 75 6D 62 65  _id=1&card_numbe
  0710: XX XX XX XX XX XX XX XX XX XX XX XX XX XX XX XX   r=4060xxxx413xxx
  0720: 39 36 26 63 61 72 64 5F 65 78 70 5F 6D 6F 6E 74  96&card_exp_mont
  0730: 68 3D 30 32 26 63 61 72 64 5F 65 78 70 5F 79 65  h=02&card_exp_ye
  0740: 61 72 3D 31 37 26 63 61 72 64 5F 63 76 6E 3D 31  ar=17&card_cvn=1
  0750: 30 39 F8 6C 1B E5 72 CA 61 4D 06 4E B3 54 BC DA  09.l..r.aM.N.T..
\end{alltt}

\begin{list2}
\item Obtained using Heartbleed proof of concepts - Gave full credit card details
\item "can XXX be exploited" - yes, clearly! PoCs ARE needed\\
without PoCs even Akamai wouldn't have repaired completely!
\item The internet was ALMOST fooled into thinking getting private keys from Heartbleed was not possible - scary indeed.
\end{list2}

\slide{sslscan}

\begin{alltt}\small
root@kali:~# sslscan --ssl2 web.kramse.dk
Version: 1.10.5-static
OpenSSL 1.0.2e-dev xx XXX xxxx

Testing SSL server web.kramse.dk on port 443
...
  SSL Certificate:
Signature Algorithm: sha256WithRSAEncryption
RSA Key Strength:    2048

Subject:  *.kramse.dk
Altnames: DNS:*.kramse.dk, DNS:kramse.dk
Issuer:   AlphaSSL CA - SHA256 - G2
\end{alltt}

Source:
Originally sslscan from http://www.titania.co.uk
 but use the version on Kali

SSLscan can check your own sites, while Qualys SSLLabs only can test from hostname



\exercise{ex:sslscan}

\exercise{ex:web-site-check}

\slide{Scanning Web Applications / Web Servers }

\hlkimage{2cm}{nikto.jpg}

\begin{quote}\footnotesize
{\bf Description}
Nikto is an Open Source (GPL) web server scanner which performs comprehensive tests against web servers for multiple items, including over 6700 potentially dangerous files/programs, checks for outdated versions of over 1250 servers, and version specific problems on over 270 servers. It also checks for server configuration items such as the presence of multiple index files, HTTP server options, and will attempt to identify installed web servers and software. Scan items and plugins are frequently updated and can be automatically updated.
\end{quote}
Source: \url{https://cirt.net/Nikto2}

\begin{list2}
\item Nikto is quite old, some would say mature tool. Always recommended to run this, cheap, quick, finds interesting links
\item Whatweb can be used to check up on a site, what framework does it use, how is it configured, where is it hosted
\item Then afterwards I recommend Burp Suite professsional \url{https://portswigger.net/burp/pro}
\item The exercise booklet also has a few extra exercises about session tokens, and SQL injection
\end{list2}

\exercise{ex:nikto-webscanner}
\exercise{ex:whatweb-scanner}

\myquestionspage

\slide{File inclusion (LFI / RFI)}
% FIXME not in these decks -- add local/remote file inclusion (obj 7)
\begin{list2}
\item \textbf{FIXME:} not in these decks -- add local/remote file inclusion (obj 7)
\end{list2}

\slide{Security misconfiguration}
% FIXME defaults, verbose errors, missing security headers (obj 7)
\begin{list2}
\item \textbf{FIXME:} defaults, verbose errors, missing security headers (obj 7)
\end{list2}

\slide{Summary}
\begin{list2}
\item Main web weaknesses: injection, XSS, CSRF, sessions, file inclusion, misconfiguration
\item Next: tools, labs and defences
\end{list2}

\end{document}
